Работа посвящена получению строгой нижней оценки константы Чватала--Санкоффа $\gamma_\sigma$ --- предела отношения математического ожидания длины наибольшей общей подпоследовательности (LCS) двух случайных равномерных строк длины $n$ над алфавитом размера $\sigma$ к $n$ при $n \to \infty$. Значение $\gamma_\sigma$ известно только в виде интервала: для $\sigma = 2$ доказано $\gamma_2 \in [0{,}7927;\, 0{,}8263]$ (Heineman et al., 2024; Lueker, 2009), при этом численные оценки показывают $\gamma_2 \approx 0{,}811$.

Предложен новый метод вычисления нижних оценок --- оконное динамическое программирование с ограниченным обзором $N$ и максимальным сдвигом $\mathrm{MAX\_SH}$. Идея доказательства: ДП вычисляет оптимальное среднее число совпадений стратегии с ограниченным обзором за $K$ шагов; по индукции и леммой Фекете о супераддитивных последовательностях нормированный предел $f_K(0,0,0)/K$ существует и не превосходит $\gamma_\sigma$, поскольку всякая такая стратегия реализуема и на реальных случайных строках. Сложность алгоритма по времени $O(K \cdot \mathrm{MAX\_SH} \cdot \sigma^{2N+2})$, по памяти $O(\mathrm{MAX\_SH} \cdot \sigma^{2N})$. Метод универсально применим для произвольного $\sigma \geq 2$. Реализация на C++ с использованием AVX2 и одинарной точности позволила достичь $N = 12$, $\mathrm{MAX\_SH} = 30$, $K = 40\,000$ при $\approx 2$~ГБ памяти и $\approx 5$~часах счёта на одном ядре.

Получены новые рекордные нижние оценки: $\gamma_2 \geq 0{,}7964$ (прирост $+0{,}48\,\%$ к Heineman 2024), а также рекорды для всех $\sigma \in \{3, \ldots, 10\}$ с приростом от $+0{,}17\,\%$ до $+5{,}42\,\%$ по сравнению с лучшими известными результатами. Все оценки независимо верифицированы методом Monte-Carlo на строках длины до $20\,000$ (по $500$~испытаний на $\sigma$).
